% Lösungen Einheit 3 von 4 – Parameter aus Maßbedingungen (zusammenbau v0.1)
\einheitenkopf{Einheit 3 von 4 · Parameter aus Maßbedingungen}

\erg{14}{a) für alle Mitglieder – der Abstandsterm muss $k$ verlieren \quad b) $\tfrac{5}{\sqrt{k^2 + 34}} = \tfrac{1}{\sqrt{2}}$, also $k^2 + 34 = 50$; $k_1 = 4$, $k_2 = -4$ \quad c) $\mathrm{sin}\,30^\circ = \tfrac{1}{\sqrt{k^2 + 2}}$, also $k^2 + 2 = 4$; $k_1 \approx 1{,}41$, $k_2 \approx -1{,}41$ \quad d) $\tfrac{|4 - k|}{3} = 2$, also $4 - k = \pm 6$; $k_1 = -2$, $k_2 = 10$ \quad e) $BD$ fest, also kleinste Höhe $MQ_k$ ($M$ Mitte von $BD$): $F_k$ senkrecht zu $SC$, $(k | k | k - 4) = r \cdot (2 | 2 | -3)$ ergibt $-3k = 2k - 8$, also $k = 1{,}6$}
\erg{15}{a) Fehler: für den Winkel Gerade–Ebene gilt der Sinus. Richtig: $\tfrac{3}{\sqrt{k^2 + 10}} = \tfrac{\sqrt{3}}{2}$, also $k^2 + 10 = 12$; $k_1 \approx 1{,}41$, $k_2 \approx -1{,}41$ \quad b) Die Ebene kann auf beiden Seiten von $P$ liegen: die Betragsgleichung lässt beide Vorzeichen zu, und jedes liefert einen Wert.}
